\documentclass[twoside, 11pt]{article}
\usepackage[utf8]{inputenc}
\usepackage[T1]{fontenc}

\usepackage[a4paper,  left=1in, right=1in, top=1in, bottom=1in, includehead, headheight=14mm]{geometry}
\usepackage[spanish]{babel}
\usepackage{graphicx}
\usepackage{subcaption}
\usepackage{amssymb}
\usepackage{ninecolors}
\usepackage{titlesec}
\usepackage[pdfusetitle]{hyperref}
\usepackage{fancyhdr}
\usepackage{lipsum}
\usepackage{tikz}
\usepackage{csquotes}

\decimalcomma

\usetikzlibrary{positioning, arrows.meta, fadings}

\hypersetup{colorlinks=true, linkcolor=black, urlcolor=black}

\pagestyle{fancy}
\fancyhf{}
\fancyhead[RE]{\nouppercase{\textsc{\leftmark}}}
\fancyhead[RO]{\nouppercase{\textsc{\rightmark}}}
\fancyhead[L]{\thepage}

\title{Microcuásares}
\author{Berredo de la Colina, Lucas \\Sánchez García, María José}

\pgfdeclareradialshading{diffusestar}{\pgfpointorigin}{%
    color(0bp)=(black!80);
    color(10bp)=(black!55);
    color(20bp)=(black!25);
    color(28bp)=(black!6);
    color(32bp)=(white)
  }

\begin{document}
\pagenumbering{roman}
\begin{titlepage}
  \centering
  {\Huge Microcuásares\par}

  \vspace{5em}

  \includegraphics[width=0.75\linewidth]{Emblema Universidad de Oviedo.jpg}\par
  \vspace{1em}
  {\huge \color{gray6} Facultad de Ciencias\par}

  \vfill

  {\Large Berredo de la Colina, Lucas\par}
  \vspace{1em}
  {\Large Sánchez García, María José\par}
  \vspace{2em}
\end{titlepage}

\thispagestyle{plain}

\noindent \textbf{Abstract.} \lipsum[1]

\vspace{2em}

\noindent \textbf{Resumen.} \lipsum[2]

\newpage

\thispagestyle{plain}
\tableofcontents
\newpage


\pagenumbering{arabic}

% \section{Introducción}
%%%%%%%%%%%%%% SUBSECCIONES %%%%%%%%%%%%%%
% Definición microcuásar                 %
% Anatomía de un sistema binario         %
%                                        %
%%%%%%%%%%%%%%%%%%%%%%%%%%%%%%%%%%%%%%%%%%
\section{Introducción}

\subsection{Definición de microcuásar y analogía de escala con los cuásares}
\noindent Un microcuásar (término acuñado por I. F. Mirabel y L. F. Rodríguez en 1992) es un sistema estelar binario compuesto por dos cuerpos ligados gravitacionalmente: el remanente compacto de una estrella masiva (siendo esta una estrella de neutrones o un agujero negro de masa estelar) y una estrella compañera de la cual extrae materia. 

Tienen su origen en el descubrimiento de chorros de radio bipolares de escala de años luz proveniendo de la fuente de rayos X duros y $\gamma$ blandos 1E 1740.7-2942, conocida como el \textit{Gran Aniquilador} (que debe su nombre a la radiación de aniquilación electrón-positrón a $511$ keV en la región del centro galáctico \cite{mirabel,mornas2026bloque1}). Estos chorros, denominados \textit{jets}, y la radiación electromagnética de alta energía son resultado directo de la actividad de los microcuasares \cite{chaty,mirabel,paredes}. 

El estudio de este tipo de sistema binario es de gran utilidad para poder entender el comportamiento de otros como pueden ser las radiogalaxias y los cuásares extragalácticos, ya que se ha demostrado a través de observaciones que replican exactamente la morfología de lóbulos y chorros características de estos \cite{mirabel}. 

Dicha unidad física existente entre los cuásares, microcuásares y colapsares progenitores de estallidos de rayos $\gamma$ largos se debe a que comparten tres características fundamentales \cite{mirabel}:
\begin{itemize}
    \item[$\bullet$] Un objeto central compacto dotado de un pozo de potencial gravitatorio muy intenso y \textit{spin}.
    \item[$\bullet$] Un disco de acreción rotante calentado por disipación viscosa y turbulenta.
    \item[$\bullet$] \textit{Jets} impulsados por procesos magnetohidrodinámicos a lo largo del eje de rotación del sistema.
\end{itemize}

En el contexto de la relatividad general, los radios gravitacional $R_g$ y de Schwarzschild $R_S$ son los que fijan la escala espacial característica de un agujero negro \cite{chaty,mornas2026bloque1}:
\begin{equation}
    R_g=\frac{GM}{c^2}\simeq 1.48\,\left(\frac{M}{M_\odot}\right)\mbox{ km}\qquad R_S=2 R_g\simeq 2.95\,\left(\frac{M}{M_\odot}\right)\mbox{ km}
\end{equation}
Como además
\begin{equation}
    t_{din}\sim\frac{R_S}{c}=\frac{2GM}{c^3}
\end{equation}
se tiene que \cite{chaty,mirabel}
\begin{equation}
    \frac{t_{MC}}{t_{AGN}}=\frac{M_{MC}}{M_{AGN}}\sim 10^{-5}-10^{-8}
\end{equation}
Por tanto, los ciclos de acomplamiento acreción-eyección que en cuásares evolucionan en escalas temporales de años a siglos se pueden observar en microcuásares en escalas de meses a un año (ciclo completo de estallido) o incluso de segundos a minutos (en fenómenos rápidos). Además, los chorros de los cuásares alcanzan longitudes de kpc o Mpc ($\sim 10^{24}$ cm), mientras que en microcuásares se extienden típicamente hasta escalas de pc ($\sim 10^{17}$ cm). Este hecho da aún más peso a la importancia del estudio de microcuásares; gracias al paralelismo latente entre estas dos estructuras tan estrechamente relacionadas podemos observar en cuestión de meses lo que en cuásares tardaría años en evolucionar \cite{chaty}. 

Por otra parte, cuando el eje del \textit{jet} de un microcuásar apunta de forma prácticamente directa hacia la línea de visión terrestre, la amplificación Doppler relativista da lugar a un microblazar, análogo estelar de los blazares extragalácticos \cite{mirabel}.

\subsection{Anatomía de un sistema binario de rayos X y de un microcuásar}

\noindent Desde el punto de vista estructural, un microcuásar constituye un subconjunto de las binarias de rayos X en las que se desarrollan chorros relativistas emisores de radiación sincrotrón no térmica \cite{paredes}. Tal y como se detalla en \cite{paredes}, alrededor de un 22\% de las binarias de rayos X catalogadas clásicamente en la galaxia por Liu et al. (2006, 2007) han sido detectadas en ondas de radio. 

La anatomía completa de un microcuásar está compuesta por cinco elementos acoplados:
\begin{enumerate}
    \item \textbf{La estrella compañera} (también denominada \textbf{donante}). Atendiendo a su tipo espectral y su masa, podemos dividir los sistemas binarios en dos grandes familias:
    \begin{enumerate}
        \item \textbf{Binarias de Rayos X de Alta Masa (HMXB)}. Se trata de o bien una estrella joven, caliente y masiva de tipo espectral O o B ($M\geq 10 M_\odot$ y $T_{eff}\sim 10^4-5\times^4$ K) con intensa emisión ultravioleta; o bien de una estrella de Wolf-Rayet (tipo espectral W, correspondiente a núcleos de He y CNO con vientos estelares fuertes \cite{mornas2026bloque2}). 
        En este tipo de estrellas, la transferencia de masa hacia el objeto compacto se produce por la captura gravitatoria del viento estelar supersónico a través de un disco de decreción ecuatorial o por el desbordamiento del lóbulo de Roche.
        \item \textbf{Binarias de Rayos X de Baja Masa (LMXB)}. En este caso, la compañera es una estrella evolucionada o de secuencia principal de baja masa y de tipo espectral tardío ($M\leq 1-2 M_\odot$ y tipo posterior a B o gigante roja de baja masa \cite{fender_gallo,malzac}). 
        Como estas estrellas carecen de vientos intensos, la transferencia de masa es únicamente consecuencia del desbordamiento del lóbulo de Roche a través del punto interno de Lagrange $L_1$.
    \end{enumerate}
    \item \textbf{El objeto compacto acretor}. Responsable de la fuerza gravitatoria del sistema. Mediante medición por espectroscopía óptica o infrarroja del desplazamiento Doppler radial de las líneas de absorción de la estrella compañera se puede determinar su masa \cite{mornas2026bloque1}. Una vez calculada, se determinan las variables necesarias a partir de las cuales se obtiene la función de masa \cite{mornas2026bloque1} dada por las leyes de Kepler:
    \begin{equation}
        f(M_1)=\frac{(M_1\sin{i})^3}{(M_1+M_2)^2}=\frac{4\pi^2(a_2\sin{i})^3}{GP_{orb}^2}=\frac{P_{orb}K_2^3}{2\pi G}
    \end{equation}
    donde $M_1$ y $M_2$ son las masas del objeto compacto y de la estrella donante respectivamente, $i$ es la inclinación del plano orbital respecto al plano del cielo, $P_{orb}$ es el periodo orbital, $K_2$ es la semiamplitud de velocidad y $a_2\sin{i}$ es el semieje proyectado. 
    
    \newpage
    De esta forma, $f(M_1)$ constituye una cota inferior absoluta para $M_1$:
    \begin{enumerate}
        \item Si $M_1\leq M_{TOV}\simeq 2-3 M_\odot$ (estabilidad de Tolman-Oppenheimer-Volkoff \cite{mornas2026bloque2}) el objeto suele ser una estrella de neutrones, donde las fuentes de tipo Z emiten intensamente en radio de forma asociada a los \textit{jets}, las fuentes \textit{Atoll} lo hacen también más débilmente. Los púlsares de rayos X fuertemente magnetizados ($B\sim 10^{12}$ G) no producen \textit{jets} de radio debido al truncamiento a gran distancia que produce la magnetosfera sobre el disco de acreción \cite{chaty}.
        \item Si $M_1>3 M_\odot$ el objeto es un agujero negro de Kerr, caracterizado de forma única por su masa $M$ y su parámetro adimensional de \textit{spin} $a_\ast=\frac{Jc}{GM^2}\in[-1,1]$ al ser eléctricamente neutro \cite{mirabel,fender_gallo}. Esto se debe a que la presión de degeneración cuántica y las fuerzas nucleares fuertes son incapaces de detener el colapso.
    \end{enumerate} 
    \item \textbf{El disco de acreción frío}. Como el momento angular específico de la materia transferida desde la estrella compañera es muy elevado, el gas se asienta en un disco de acreción kepleriano en lugar de caer radialmente hacia el objeto compacto. Este disco es geométricamente delgado ($\frac{H}{r} \ll 1$) y ópticamente grueso ($\tau \gg 1$) \cite{malzac}. El momento angular es transportado hacia el exterior por las tensiones viscosas, que disipan energía potencial gravitaria provocando un calentamiento de las regiones internas del disco hasta $T\sim 10^6-10^7$ K. De esta forma, emite rayos X blandos ($\sim0,1-2$ keV) \cite{fender_gallo,malzac}.
    \item \textbf{La corona de plasma caliente} (flujo interno de acreción caliente). La emisión de cuerpo negro del disco es incapaz de explicar la emisión de fotones de rayos X duros y rayos $\gamma$ blandos que se detecta observacionalmente. Esto es debido a que la componente dura se origina en la corona, una región ópticamente delgada de electrones altamente energéticos ($T_e\sim10^9$ K $\approx100$ keV) \cite{malzac}.
    \item \textbf{\textit{Jets} y vientos del disco}. El sistema lanza \textit{jets} perpendicularmente al plano del disco de acreción, a lo largo del eje de rotación. Estos son observados en ondas de radio, milimétricas, infrarrojas y ópticas mediante radiación de sincrotrón y en rayos X y $\gamma$ mediante dispersión Compton inversa \cite{paredes,fender_gallo,malzac}. Recientemente, también se ha descubierto que el disco expulsa vientos ecuatoriales densos altamente ionizados no relativistas durante los estados espectrales blandos en los que el \textit{jet} polar se apaga \cite{fender_gallo}.
\end{enumerate}


% \section{Línea temporal de vida}
%%%%%%%%%%%%%% SUBSECCIONES %%%%%%%%%%%%%%
%                                        %
% Requisitos para formación              %
% Origen mediante supernova              %
% Sistema binario durmiente              %
% Etapa activa y dependencias            %
% Desmantelamiento sistémico             %
%                                        %
%%%%%%%%%%%%%%%%%%%%%%%%%%%%%%%%%%%%%%%%%%
  \section{Línea temporal de vida}

\noindent El objetivo de esta sección será describir a grandes rasgos la línea temporal que habita un microcuásar. Partiendo de los requisitos, su origen, los estados durmientes y activos, y su fin. En particular, se dejará la mayor parte del estado activo a la sección siguiente, <<El ciclo de estados espectrales>>. 

Debido al flujo de la literatura, esta sección será más genérica a sistemas binarios de rayos X, sistemas binarios con agujeros negros, y centralmente, LMXB (Binarios de rayos X de baja masa, \textit{Low Mass X-ray Binaries}). El estudio de estos cuerpos es una superclase al estudio de los microcuásarses, y por tanto no deja de ser relevante.

\begin{figure}[htbp]
\centering
\begin{tikzpicture}[
  node distance=0.6cm and 0.5cm,
  every node/.style={
    draw, 
    font=\tiny, 
    rounded corners=2pt, 
    inner sep=3pt, 
    align=center,
    minimum height=11pt,
  },
  arrow/.style={-Stealth}
]

\node (node1) [text width=100pt] {Requisitos - Sistema \\binario masivo};

\node (node2) [below=of node1, text width=100pt] {Origen - Supernova};

\node (node3) [below right=10pt and -20pt of node2, text width=100pt] {Colapso del núcleo};

\node (node4) [below right=30pt and -20pt of node2, text width=100pt] {Supervivencia};

\node (node5) [below=50pt of node2, text width=100pt] {Sistema binario durmiente};

\node (node6) [below=of node5, text width=100pt] {Microcuásar activo};

\node (node7) [below right=10pt and -20pt of node6, text width=100pt] {Roche-Lobe / Vientos solares};

\node (node8) [below right=30pt and -20pt of node6, text width=100pt] {Diagrama Intensidad-Dureza};

\node (node9) [below=50pt of node6, text width=100pt] {Etapas finales};

\node (node10) [below right=10pt and -20pt of node9, text width=100pt] {Fin de \textit{jets} relativistas};

\node (node11) [below right=30pt and -20pt of node9, text width=100pt] {Muerte de estrella};

\node (node12) [below right=50pt and -20pt of node9, text width=100pt] {Unifiación del sistema};

\draw [arrow] (node1.south) -- (node2.north);

\draw [arrow, color=blue!65!black] (node2.south) |- (node3.west);

\draw [arrow, color=blue!65!black] (node2.south) |- (node4.west);

\draw [arrow] (node2.south) -- (node5.north);

\draw [arrow] (node5.south) -- (node6.north);

\draw [arrow, color=blue!65!black] (node6.south) |- (node7.west);

\draw [arrow, color=blue!65!black] (node6.south) |- (node8.west);

\draw [arrow] (node6.south) -- (node9.north);

\draw [arrow, color=blue!65!black] (node9.south) |- (node10.west);

\draw [arrow, color=blue!65!black] (node9.south) |- (node11.west);

\draw [arrow, color=blue!65!black] (node9.south) |- (node12.west);

\end{tikzpicture}
\caption{Línea del tiempo representando la vida de un microcuásar.}
\end{figure}

\subsection{Requisitos para formación. Origen.}

\paragraph{Evolución de estrellas masivas y límite de Eddington.} La formación del objeto compacto, motor del microcuásar, requiere que la estrella primaria nazca en la ZAMS con una masa suficiente tal que entre en el ciclo CNO, i.e.,

\begin{equation}
  M\geq 8-10 M_{\odot}
\end{equation}

\noindent De acuerdo con las ecuaciones de homología, las estrellas masivas tienen un núcleo convectivo (debido a que $\varepsilon_{\mbox{CNO}}\propto\rho T^{12}$, que produce una luminosidad central demasiado alta para ser evacuada únicamente por radiación, violando el Criterio de Estabilidad de Schwarzschild). La envoltura, en cambio, sí es radiativa, dominada por la dispersión de Thomson en los electrones libres. De esta forma, la relación masa-luminosidad es muy elevada ($L\propto M^3$), y la estrella vive poco tiempo en la ZAMS.

\begin{equation}
  \tau_{\mbox{nuc}}\sim \frac{\varepsilon M c^2}{L}\sim 10^6-10^7\mbox{ años}.
\end{equation}

\noindent Estas estrellas, localizadas en la región de presión de radiación en el diagrama H-R, siguen la ecuación de estado

\begin{equation}
  P_{\mbox{rad}}=\tfrac{1}{3}aT^4=(1-\beta)P.
\end{equation}

\noindent Debido a su masa, el índice adiabático tenderá hacia $4/3$, el umbral de inestabilidad. Así, es necesario imponer que la fuerza de atracción sea mayor que el gradiente de la presión de radiación. Se define la luminosidad de Eddington

\begin{equation}
  L_{\mbox{Edd}}=\frac{4\pi GMc}{\kappa}\approx\frac{4\pi GMm_pc}{\sigma_T}\approx 1,3\times10^{38}\left(\tfrac{M}{M_{\odot}}\right)\approx 3,2\times 10^4\left(\tfrac{M}{M_{\odot}}\right)L_{\odot}.
\end{equation}

\noindent Esta magnitud acabará correspondiéndose con la unidad natural de luminosidad para microcuásares. 

Para que el sistema binario acabe convirtiéndose en un microcuásar, la separación orbital debe de ser suficientemente estrecha para que haya transmisión de masa desde la estrella al objeto compacto, pero sin que se fusionen al completo. 

En el caso de que el binario tenga masas muy dispares (de al menos 25 órdenes de magnitud), se deberá cumplir una fase de envoltura común: la primaria en fase de supergigante engloba a la secundaria, y la fricción hidrodinámica extrae energía orbital. La secundaria se acaba expulsando de la envolvente, y dando lugar un sistema binario compacto, con un núcleo de helio (estrella de Wolf-Rayet), y una estrella de baja masa.

Una vez agotado el hidrógeno en el centro de la primaria, se evita el \textit{flash} de helio mediante una emisión térmica de neutrinos. El núcleo encadena las siguientes reacciones:

\vspace{-0.5em}

\begin{enumerate}
  \item Combustión nuclear de $\phantom{}^4\mbox{He}$ (reacción triple-$\alpha$, $T\geq 10^8\mbox{K}$)\vspace{-0.5em}
  \item Combustión nuclear de $\phantom{}^{12}\mbox{C}$ ($T\in[5\times10^8, 10^9]K$, para una estrella de $25M_{\odot}$, $\sim600$ años)\vspace{-0.75em}
  \item Combustión nuclear de $\phantom{}^{20}\mbox{Ne}$ ($\sim 1$ año)\vspace{-0.5em}
  \item Combustión nuclear de $\phantom{}^{16}\mbox{O}$ ($T\geq1,4\times10^9\mbox{K}$, $\sim 0,5$ años)\vspace{-0.5em}
  \item Combustión nuclear de $\phantom{}^{28}\mbox{Si}$ ($T\geq 2-4\times10^{9}\mbox{K}$, $\sim 1$ día)\vspace{-0.5em}
\end{enumerate}

\noindent Al llegar a los metales férreos ($\mbox{Fe, Ni}$), con energía de ligadura muy alta ($B/A\approx8,8\mbox{MeV}$), las reacciones dejan de ser exotérmicas, y se empieza a consumir energía térmica.

Consecuentemente, el núcleo de $\phantom{}^{56}\mbox{Fe}$, originalmente sostenido por la presión de los electrones degenerados ultrarrelativistas ($\Gamma=4/3$), sigue creciendo hasta el límite de Chandrasekhar

\begin{equation}
  M_{\mbox{Ch}}=4\pi \left[\tfrac{K}{\pi G}\right]^{3/2}=5,83\mu_e^{-2}M_{\odot}\approx 1,46 M_{\odot}
\end{equation}

\noindent En este instante se produce un colapso gravitatorio. El proceso de caída libre tiene una duración aproximada de $\tau_{\mbox{libre}}\sim100\mbox{ms}$. Esta desintegración es apoyada por dos fenómenos:

\begin{enumerate}
  \item \textbf{Fotodesintegración:} $\gamma+\phantom{}^{56}\mbox{Fe}\to13\phantom{}^4\mbox{He}+4n-124,4\mbox{MeV}$, $\gamma+\phantom{}^{4}\mbox{He}+2p+2n-28,3\mbox{MeV}$. Resultado: $\Gamma\lneq 4/3$.
  \item \textbf{Captura electrónica (neutronización):} $p+e^-\to n+\nu_e$. Resultado: liberación de neutrinos, llevando potencial gravitatorio.
\end{enumerate}

\noindent Cómo finaliza este colapso indica la naturaleza del microcuásar.

\paragraph{Caso 1. \textit{Natal kick}.} Cuando la masa de las progenitoras a edad $0$ está entre $8M_{\odot}$ y $40M_{\odot}$ y $\rho\sim 2-3\rho_0$, los nucleones forman un líquido de Fermi. Éste queda gobernado por la ecuación de Tolman-Oppenheimer-Volkoff:

\begin{equation}
  \tfrac{dP}{dr}=-\frac{Gm(r)\rho(r)}{r^2}\left[1+\frac{P(r)}{\rho(r)c^2}\right]\left[1+\frac{4\pi r^3 P(r)}{m(r)c^2}\right]\left[1-2\frac{Gm(r)}{rc^2}\right]^{-1}.
\end{equation}

\noindent Las capas externas rebotan sobre el núcleo rígido, generando una onda de choque que expulsa a las exteriores, es decir, una supernova de colapso del núcleo. (En efecto, será de tipo II si conservaba $\mbox{H}$, ó Ib/Ic si no). Queda restante una estrella de neutrones $M\sim2M_{\odot}$ o un agujero negro de baja masa. Este último caso se llama \textit{fallback}, y ocurre cuando la caída del material supera el límite de TOV.

Apelando a la <<Arqueología Galáctica>> \cite{mirabel}, la eyección de masa mediante el efecto Blaauw y acumulación de asimetrías hidrodinámicas conllevan un fuerte impulso al centro de masas, conocido como \textit{Natal kick}. 


\paragraph{Caso 2. Colapso directo.} Cuando la masa en la ZAMS supera $25-40M_{\odot}$ y se tiene una baja metalicidad, el núcleo de hierro supera directamente el límite de estabilidad de TOV. La onda de choque fracasa y la estrella implosiona sin producir una supernova luminosa ni expulsar una cantidad de masa significativa.

Los sistemas binarios resultantes tienen velocidades extremadamante bajas, normalmente asociados a rotación galáctica. Sus agujeros negros fueron formados por colapso directo. 

Si la estrella tuviera momento angular previo, se formarán tanto un agujero negro como un disco de acreción ultradenso, expulsando chorros ultrarrelativistas que perforan la estrella primaria. Ésta es perforada, y se produce un estallido de rayos gamma de larga duración.


\subsection{Sistema binario durmiente}

\noindent Si el sistema binario se ha formado de manera correcta y no colapsa, se puede permanecer una cantidad de tiempo indeterminada (y posiblemente muy larga) en un estado durmiente (también conocido como estado de quiesencia). En el caso particular de las binarias de rayos X de rayos X (LMXB), esta será la etapa más larga de su vida útil. 

Los episodios de alta intensidad (\textit{outbursts}) duran unos meses, y van intercalados por intervalos de recurrencia de años o décadas. En este periodo, el ritmo de transferencia es extremadamente bajo, 

\begin{equation}
  L_X\sim 10^{-8,5}-10^{-6,5}L_{\mbox{Edd}}.
\end{equation}

\noindent Según Malzac \cite{malzac}, el disco de acreción se halla a una gran distancia del agujero negro, mientras que la región interna contiene un flujo caliente advectivo tenue e ineficiente. (ADAF, RIAF).

Para activar un microcuásar durmiente, uno debe referirse al modelo de Meyer-Hofmeister, basado en la inestabilidad térmica-viscosa asociada a la ionización del hidrógeno en el disco de acreción. Los pasos a seguir son:

\begin{enumerate}
  \item El gas transferido desde la secundaria hasta la primaria se acumula en las regiones externas del disco. El hidrógeno está en el estado neutro de la ecuación de Saha. La viscosidad efectiva es minúscula, el flujo entrante sobrepasa al cebado del núcleo.
  \item Según la masa se acumula en el exterior, crece la densidad y temperatura. Entorno a los $6000$ ó $10000$ Kelvin, se ioniza el hidrógeno.
   
  En este régimen, la opacidad de Rosseland aumenta (debido a $\mbox{H}^-$ y las transiciones \textit{bb}). Además, el índice adiabático ($\Gamma=\partial_{\ln\rho}\ln P$) cae: $4/3\to 1\mbox{ o }2$, provocando inestabilidad convectiva.
  \item Cuando $T\geq 10^4\mbox{K}$, el gas está totalmente ionizado. El acoplamiento magnético del cono de plasma dispara la viscosidad. Se desencadena una avalancha de materia, aumentando la tasa de acreción y la luminosidad. 
\end{enumerate}


\begin{figure}[h]
\centering
\makebox[\textwidth][c]{
\begin{tikzpicture}[>=Stealth, font=\sffamily]

\begin{scope}[xshift=0cm]
    \node[above=1.4cm, font=\bfseries] at (-0.3,0) {Binaria primordial};
    
    \draw[dash pattern=on 2.5pt off 2pt, line width=0.6pt] 
      (-0.45, 0) to[out=135, in=0] (-1.2, 0.6)
      arc (90:270:0.6)
      to[out=0, in=225] (-0.45, 0);

    \draw[dash pattern=on 2.5pt off 2pt, line width=0.6pt] 
      (-0.45, 0) to[out=45, in=180] (0.5, 1.2)
      arc (90:-90:1.2)
      to[out=180, in=-45] (-0.45, 0);

    \fill (-1.2, 0) circle (2.2pt);
    \fill (0.5, 0) circle (6.5pt);
  \end{scope}

  \draw[->, line width=0.7pt] (1.9, 0) -- (2.6, 0);

  \begin{scope}[xshift=4.6cm]
    \node[above=1.4cm, font=\bfseries] at (-0.3,0) {Roche-lobe};
    
    \fill[shading=diffusestar] (0.5, 0) circle (1.1);

    \draw[dash pattern=on 2.5pt off 2pt, line width=0.6pt] 
      (-0.45, 0) to[out=135, in=0] (-1.2, 0.6)
      arc (90:270:0.6)
      to[out=0, in=225] (-0.45, 0);

    \draw[dash pattern=on 2.5pt off 2pt, line width=0.6pt] 
      (-0.45, 0) to[out=45, in=180] (0.5, 1.2)
      arc (90:-90:1.2)
      to[out=180, in=-45] (-0.45, 0);

    \fill (-1.2, 0) ellipse (0.38 and 0.08);
    \fill (-1.2, 0) circle (1.8pt);
    \draw[line width=0.7pt] (-0.45, 0) to[out=180, in=0] (-0.85, -0.02);
    \fill (0.5, 0) circle (6.5pt);
  \end{scope}

  \draw[->, line width=0.7pt] (6.7, 0) -- (7.4, 0);

  \begin{scope}[xshift=9.2cm]
    \node[above=1.4cm, font=\bfseries] at (0,0) {Espiral};
    
    \fill[shading=diffusestar] (0, 0) circle (1.15);
    \fill (0, 0) circle (7pt);
    \draw[line width=0.6pt] (0, 0) ellipse (0.85 and 0.8);
    \fill (-0.5, 0.65) circle (2.6pt);
  \end{scope}

  \draw[->, line width=0.7pt] (10.65, 0.35) -- (11.35, 0.85);
  \draw[->, line width=0.7pt] (10.65, -0.35) -- (11.35, -0.85);

  \begin{scope}[xshift=11.9cm, yshift=0.9cm]
    \begin{scope}[scale=0.6]
      \draw[line width=0.6pt] 
        (-0.25, 0) to[out=135, in=0] (-0.7, 0.35)
        arc (90:270:0.35)
        to[out=0, in=225] (-0.25, 0);

      \draw[line width=0.6pt] 
        (-0.25, 0) to[out=45, in=180] (0.4, 0.65)
        arc (90:-90:0.65)
        to[out=180, in=-45] (-0.25, 0);

      \fill (-0.7, 0) circle (2.2pt);
      \fill (0.4, 0) circle (4.5pt);
    \end{scope}
    \node[anchor=west, font=\small\bfseries] at (0.7, 0) {Binario compacto};
  \end{scope}

  \begin{scope}[xshift=11.9cm, yshift=-0.9cm]
    \fill (0.1, 0) circle (3pt);
    \node[anchor=west, font=\small\bfseries] at (0.7, 0) {Estrella aislada};
  \end{scope}

\end{tikzpicture}}
\caption{Recreado desde XiangDong 2010.}
\end{figure}

\vspace{1em}

\subsection{Etapa activa}

\noindent Se refiere al lector a la sección siguiente.



\subsection{Desmantelamiento sistémico}

\noindent La etapa activa concluye cuando el suministro de materia de la estrella secundaria acaba. La estructura disco-corona-\textit{jet} se desmantela sistémicamente. Existen dos mecanismos posibles para esto, dependiendo de la masa restante.

\begin{enumerate}
  \item \textbf{LMXB:} La estrella donante pierde su envoltura de $\mbox{H}$. El núcleo remanante se enfría como una enana blanca, dejando un sistema binario pasivo: una enana blanca y un agujero negro/púlsar.
  \item \textbf{HMXB:} En binarias de alta masa (p.ej. Cygnus X-1, Cygnus X-3), la estrella donante (O/B/Wolf-Rayet) agota su combustible nulcear y se colapsa el núcleo por la fuerza gravitatoria. Si la pérdida de masa o el segundo \textit{Natal Kick} (esta vez asimétrico) supera la energía de ligadura, el sistema presenta una disrupción. Ambos compactos son disparados al espacio interestelar, en distintos caminos con altas velocidades. Si no se desligan, se forma un binario de dos objetos compactos: NS-NS (PSR 1913+16), púlsar doble (J0737-3039), NS+BH ó BH+BH.
\end{enumerate}

\noindent A partir de este momento, la evolución final es gobernada por la relatividad general. El momento cuadripolar de ambas masas compactas,

\begin{equation}
  Q_{ij}=\int\rho(x)(x_ix_j-\tfrac{1}{3}r^2\delta_{ij})dx
\end{equation}

\noindent da lugar a ondas gravitacionales. Su potencia cuadripolar es

\begin{equation}
  L_{GW}=\tfrac{1}{5}\tfrac{G}{c^5}\left\langle \dddot{Q}_{ij}, \dddot{Q}^{ij}\right\rangle\sim \frac{c^5s^2}{G}\left(\frac{R_S}{R}\right)^2\left(\frac{v}{c}\right)^6.
\end{equation}

\noindent La pérdida de energía conlleva una reducción en el radio orbitial. El periodo orbital estará gobernado por

\begin{equation}
  \Delta T(t)=\frac{\dot{P}_{\mbox{orb}}}{2P_{\mbox{orb}}}t^2
\end{equation}

\noindent en un tiempo $\tau\sim10^7-10^9$ años. 

Las fases finales (\textit{inspiral}, \textit{merger}, \textit{ringdown}) son detectadas en la Tierra por interferómetros como LIGO.


% Véase ACCRETION TURNOFF AND RAPID EVAPORATION OF VERY LIGHT SECONDARIES IN LOW-MASS X-RAY BINARIES M. Ruderman, J. Shaham, and M. Tavani





% \section{Etapa activa. El ciclo de estados espectrales}
\section{El ciclo de estados espectrales}
%%%%%%%%%%%%%% SUBSECCIONES %%%%%%%%%%%%%%
% Diagrama HID                           %
%                                        %
%%%%%%%%%%%%%%%%%%%%%%%%%%%%%%%%%%%%%%%%%%

% \section{Experimentos numéricos}
\section{Experimentos numéricos}

\noindent El objetivo de esta sección es realizar una serie de diagramas y resultados numéricos, partiendo de simulaciones o datos experimentales.

\subsection{Recreación Diagrama Intensidad-Dureza}

% Descripción matemática del diagrama
% Fuentes de datos experimentales
% Implementación
% Resultados

\subsection{Cinemática de los \textit{jets}. Movimiento supralineal aparente.}

% Descricpión matemática del problema
% Simulación
% Resultados

\subsection{Difusión viscosa del disco de acreción}

% Descripción matemática del problema
% Método de resolución numérica de EDP
% Resultados


% \section{Conclusiones}
\section{Conclusiones}
\noindent Los microcuásares y los cuásares manifiestan representan, a diferentes escalas, el mismo tipo de fenómenos físicos. Dado que su radio gravitacional y escala temporal dinámica escalan linealmente con la masa del agujero negro, los microcuásares sufren los mismos procesos que a los cuásares les tomarían siglos, e incluso miles de años, en escalas humanas.  \\

\noindent El estudio cinemático y evolutivo de los microcuásares permite reconstruir cómo mueren las estrellas masivas, demostrando que la masa del remanente determina si el colapso ocurre de forma explosiva o silenciosa. Tras su fase de alta energía, estos sistemas culminan su ciclo como binarias compactas destinadas a fusionarse mediante la emisión de ondas gravitacionales. \\

\noindent Los estados espectrales de los microcuásares reflejan distintas geometrías de acreción: un disco delgado que llega hasta el $R_{\text{ISCO}}$ en el estado blando, frente a un estado duro dominado por Comptonización térmica. Para explicar este último se plantean dos modelos principales: un disco frío truncado con un flujo interior caliente (ADAF/RIAF), o una corona dinámica eyectada (\textit{outflowing} ADC) vinculada a la base del \textit{jet}. \\

\noindent Los estallidos de las binarias de rayos X describen un ciclo de histéresis en forma de <<Q>> gobernado por la asimetría entre la evaporación del disco (estado duro) y el enfriamiento Compton (estado blando). Esto impone una alternancia excluyente en la eyección de materia: el estado duro emite jets polares relativistas (divididos en dos pistas según su eficiencia radiativa), mientras que el estado blando los suprime y activa vientos ecuatoriales ionizados. Al transicionar entre ambos estados, el colapso del disco hasta la última órbita estable ($R_{\text{ISCO}}$) apaga el \textit{jet} continuo y lanza eyecciones discretas ultrarrelativistas. 

% --- BIBLIOGRAFÍA ---
\cleardoublepage
\begin{thebibliography}{99}

\bibitem{chaty}
Chaty, S. 2022. \textit{Accreting Binaries}, cap.~8: ``Jets in accreting binaries''. IoP Publishing.

\bibitem{mirabel}
Mirabel, F. 2010. ``Microquasars: summary and outlook'', en \textit{The Jet Paradigm: From Microquasars to Quasars}, ed. T. Belloni, \textit{Lecture Notes in Physics}, Vol.~794. Springer-Verlag, págs.~1--15.

\bibitem{paredes}
Paredes, J.~M. 2011. ``Radio jets and high energy emission in microquasars'', \textit{Mem. Soc. Astron. Ital.}, 82, 174--181.

\bibitem{fender_gallo}
Fender, R. \& Gallo, E. 2014. ``An Overview of Jets and Outflows in Stellar Mass Black Holes'', \textit{Space Sci. Rev.}, 183, 323--337.

\bibitem{malzac}
Malzac, J. 2007. ``Models for microquasars'', arXiv:0706.2389.

\bibitem{mornas2026bloque1}
Mornas, L. \& Toffolatti Ballarin, L. 2026. \textit{Astrofísica y Cosmología: Bloque 1 -- Lo observable en astrofísica} (Diapositivas de clase). Universidad de Oviedo.

\bibitem{mornas2026bloque2}
Mornas, L. 2026. \textit{Astrofísica y Cosmología: Bloque 2 -- Física de las estrellas} (Diapositivas de clase). Universidad de Oviedo.

\end{thebibliography}



\end{document}
